\documentclass[12pt]{article}
\usepackage{graphicx}
\usepackage{hyperref}
\usepackage{amsmath}
\usepackage{amssymb}
\usepackage{siunitx}
\usepackage[numbers, square, sort&compress]{natbib}

\title{Chapter 6: Preparation of Small-Molecule Libraries}
\author{Elvis Martis}

\begin{document}
\maketitle

\section{Introduction}

Small molecules that bind receptors are called ligands. Although the terms \emph{ligand} and \emph{drug} are sometimes used interchangeably, it is important to recognise that while all drugs are ligands, not all ligands qualify as drugs. Ligands are typically represented initially as 2D atomic connectivity graphs. These representations are converted into 3D structures and subjected to energy minimisation before molecular docking simulations. When screening large chemical libraries, often comprising millions or even billions of compounds, an initial filtering stage is commonly applied to enrich for molecules with drug-like properties. This process typically involves evaluating physicochemical descriptors such as molecular weight, lipophilicity (log P), the number of rotatable bonds, and the counts of hydrogen bond donors and acceptors. One of the most widely applied guidelines for such filtering is Lipinski’s Rule of Five, originally proposed by Christopher Lipinski\cite{CH6_lipinski1997}. It should be noted that the Rule of Five primarily applies to orally administered drugs that permeate biological membranes through passive diffusion rather than active or receptor-mediated transport mechanisms.
Nevertheless, these filters remain valuable for excluding compounds that are unlikely to exhibit suitable pharmacokinetic properties. In addition to physicochemical filtering, the identification and removal of pan-assay interference compounds (PAINS)\cite{CH6_baell2010} constitute an important step in early screening workflows. PAINS are chemical structures that frequently appear as active across multiple assay systems and are therefore often associated with assay interference or false-positive activity. 

In most docking campaigns, the computational protocol is only as reliable as the quality of the input receptor structure, and equally critically, the representation and preparation of small molecules. Inadequate ligand preparation, such as incorrect protonation, missing tautomers, unrealistic conformations, or inconsistent valence and stereochemistry, can lead to misleading enrichments, even if the underlying search and scoring algorithms are otherwise robust.\cite{CH6_Sastry2013,CH6_kolb2009,CH6_KeyTopicsDocking} 

This chapter aims to provide a systematic and in-depth overview of methods and software for preparing and handling small molecules for molecular docking. The focus is on both theoretical concepts (e.g., protonation equilibria, tautomerism, conformational sampling, charge models) and practical aspects (software packages, file formats, and databases). Emphasis is placed on reproducible workflows, verification of structural correctness, and the use of curated ligand datasets and benchmarks. 

Ligand preparation, therefore, aims to:
\begin{enumerate}
  \item Standardise connectivity and valence to ensure correct bond orders, aromaticity, and formal charges consistent with chemical intuition and experimental p$K_a$/redox state.
  \item Generate relevant protonation and tautomeric states to enumerate or select those states likely populated at the target pH and environment.
  \item Resolve stereochemistry and isomerism to appropriately define chiral centres and E/Z double-bond stereochemistry, enumerating alternatives if necessary.
  \item Generate low-energy 3D conformations by sampling accessible conformational space so that docking can identify plausible binding modes.
  \item Assign force field parameters and partial charges to provide optimal parameters required by the docking engine's scoring function.
  \item Filter and quality-control to remove impossible or highly strained structures, salts, and reactive or toxic chemotypes.
\end{enumerate}

Extensive benchmarking has shown that systematic attention to these steps can markedly improve virtual screening enrichment factors, sometimes by a factor of 2 or more\cite{CH6_Sastry2013,CH6_DEKOIS2}.

\subsection{Typical workflow}

A typical ligand preparation workflow for docking can be outlined as:

\begin{enumerate}
  \item Input 1D/2D structures (SMILES or SDF) and perform basic corrections, such as valence checks, and standardisation.
  \item Delete salts and solvents atoms.
  \item Predict and enumerate relevant protonation states and tautomers at specified pH (e.g., 7.4 \(\pm\) 1).
  \item Enumerate undefined stereocenters and E/Z isomers, within user-defined limits.
  \item Generate 3D coordinates for each state and perform conformational sampling.
  \item Assign partial charges and any special atom types needed by the docking program.
  \item Perform restrained energy minimisation using a force field suitable for small molecules.
  \item Export to the docking engine's preferred file format (e.g., MOL2, PDBQT, or SDF).
  \item Verify geometry and chemistry with automated and manual checks.
\end{enumerate}

Software such as Open Babel\cite{CH6_OBoyle2011}, RDKit (with ETKDG-based embedding)\cite{CH6_Riniker2015}, OMEGA\cite{CH6_Hawkins2010}, and Gypsum-DL\cite{CH6_Ropp2019} implement many of these steps in open-source form, while commercial suites (Schr\"odinger, OpenEye, MOE, Discovery Studio) provide additional proprietary functionality.

\section{Representation and Basic Processing of Small Molecules}

\subsection{1D, 2D, and 3D representations}

Small molecules may be initially represented as:
\begin{itemize}
  \item 1D line notations: e.g., SMILES\cite{CH6_Weininger1988}, InChI\cite{CH6_Heller2013}. These encode molecular graphs as strings, which are convenient for storage in large databases and for substructure/similarity searches.
  \item 2D connection tables: e.g., MDL MOL/SDF, which store atom types, bond orders, and 2D coordinates. Molecular formats such as SDFs allow the storage of molecular properties.
  \item 3D coordinate files: e.g., MOL2, PDB, SDF or PDBQT, which include 3D coordinates and partial charge information suitable for docking. The amount of information needed depends on the docking software one chooses.
\end{itemize}

Docking requires 3D geometries; however, ligand databases, such as ChEMBL\cite{CH6_Gaulton2012}, ZINC\cite{CH6_Irwin2005}, and PubChem\cite{CH6_KimPubChem}, often provide structural information in 1D SMILES or 2D SDF formats, necessitating the generation of 3D coordinates and conformers.

\subsection{Sanitisation}

Most cheminformatics toolkits implement a sanitisation step that verifies chemical correctness. These include valence checks to ensure that valence counts of atoms conform to allowed patterns (e.g., tetravalent carbon, trivalent nitrogen in amines). Furthermore, valence checks also aid in correcting missing bond orders from connectivity. Since most drugs are aromatic in nature, a special emphasis on aromaticity is crucial, and applying models such as Huckel-like rules or pre-tabulated aromatic atom types is helpful for labelling aromatic rings. Lastly, the partial assignment to ensure the total charge is consistent with valence and, where possible, with expected p$\textrm{K}_\textrm{a}$.

In RDKit\cite{CH6_bento2020open}, for example, sanitisation is performed via \verb|Chem.SanitizeMol|, which checks valence, aromaticity, conjugation, and ring information, raising errors or warnings when invalid structures are encountered. Open Babel\cite{CH6_OBoyle2011} provides analogous functionality, including detection of impossible valence and inconsistent aromaticity. Sanitisation should be performed \emph{before} any optimisation or docking-related processing. 

\subsection{Standardization}
Standardisation aims to convert diverse input representations into a consistent internal representation. Common steps include:
\begin{itemize}
  \item Deleting salts to remove counterions (e.g., chloride, sodium) and solvents (e.g., water, DMSO) using fragment-based rules.
  \item Normalisation of functional groups to convert alternative representations (e.g., nitro as $\textrm{N}^{+}(\textrm{=O})[O^-]$) into a canonical pattern.
  \item Neutralisation to adjust charges on strong acids/bases to a standard reference pH or to a neutral form.
  \item Stereochemistry cleanup to remove or set undefined chiral centres when the input stereo is ambiguous.
\end{itemize}

Standardisation libraries, such as e.g., RDKit's \verb|rdMolStandardize| and commercial standardisation tools, encapsulate rules derived from medicinal chemistry practice.

\section{Protonation States, Tautomerism, and Stereochemistry}

\subsection{Protonation equilibria and the Henderson--Hasselbalch equation}

The protonation state of ionizable groups, such as carboxylic acids and amines, is governed by acid-base equilibria. 
For a monoprotic acid HA,

\begin{equation}
\textrm{HA} \rightleftharpoons \textrm{A}^- + \textrm{H}^+
\end{equation}

The Henderson--Hasselbalch equation gives the ratio of deprotonated to protonated species at equilibrium:
\begin{equation}
\textrm{pH} = \textrm{p}K_a + \log_{10} \frac{[\textrm{A}^-]}{[\textrm{HA}]} 
\end{equation}
Rearranging,
\begin{equation}
\frac{[\textrm{A}^-]}{[\textrm{HA}]} = 10^{\textrm{pH} - \textrm{p}K_a}
\end{equation}

For a given pH and known p$\textrm{K}_\textrm{a}$, this allows the computation of the fractional populations of each protonation microstate. For multiple ionizable centres, microstate populations can be computed from a micro-$\textrm{K}_\textrm{a}$ or micro-p$\textrm{K}_\textrm{a}$ model; in practice, software such as Epik\cite{CH6_shelley2007epik}, ChemAxon p$\textrm{K}_\textrm{a}$ predictors, or Dimorphite-DL\cite{CH6_ropp2019dimorphite} (used within Gypsum-DL\cite{CH6_Ropp2019}) implement empirical or semi-empirical models to enumerate and score protonation microstates\cite{CH6_Sastry2013}.

In docking, it is rarely sufficient to consider a single protonation state; instead, all states within a reasonable free energy window (often a few kcal/mol at the working pH) should be generated. The Boltzmann weight of each state i is

\begin{equation}
p_i = \frac{e^{-\beta G_i}}{\sum_j e^{-\beta G_j}}, \quad \beta = \frac{1}{k_B T}
\end{equation}

where $\textrm{G}_\textrm{i}$ is the free energy of microstate i, and $\textrm{k}_\textrm{B}$ is the Boltzmann contant. At a typical screening pH, states whose populations are below 1--5\% may be neglected; however, borderline cases, such as histidines and anilines, can strongly influence binding modes.

\subsection{Tautomerism}

Many heterocycles and functional groups exhibit tautomerism, such as keto-enol, amide-imidic acid, and lactam-lactim. Tautomers often differ in hydrogen placement and formal charges, with potentially large effects on hydrogen bonding and electrostatics in the binding site. Automatic tautomer enumeration proceeds by defining substructure-based transformation rules (SMARTS patterns) that map between tautomeric forms.

In practice:
\begin{itemize}
  \item Enumerators generate all possible tautomers under a given set of rules.
  \item Scoring functions or empirical rules estimate relative stability (e.g., favouring amide over imidic acid forms).
  \item A subset of low-energy tautomers is selected for docking, sometimes weighted by estimated populations.
\end{itemize}

Tools such as Gypsum-DL explicitly enumerate tautomers (and other isomeric forms) and produce 3D structures for each selected form\cite{CH6_Ropp2019}. Commercial programs (e.g., LigPrep, Epik, and MOE) implement proprietary tautomer generators and scoring schemes\cite{CH6_Sastry2013}.

\subsection{Stereochemistry and E/Z isomerism}

Stereochemical ambiguities are common in large virtual libraries. The chiral centres may be unspecified (e.g., input from generic SMILES), while the E/Z configuration of alkenes may be unknown or partially constrained. Atropisomerism and conformational chirality are usually not explicitly encoded. A general strategy is to keep experimentally determined stereochemistry fixed from curated databases such as ChEMBL or PDB ligands. This is followed by enumerating all reasonable stereoisomers for undefined centres; however, cap the total number of isomers per molecule to control combinatorial explosion, and for flexible or unassigned stereocenters, consider whether stereochemistry is relevant to the docking hypothesis. Gypsum-DL\cite{CH6_Ropp2019}, for example, can enumerate R/S stereocenters, E/Z isomers, and ring conformers, within user-defined limits, for each input molecule.

\section{3D Coordinate Generation and Conformational Sampling}
The goal of conformer generation is to produce an ensemble of low-energy 3D geometries that approximate the accessible conformational space of a molecule in solution, while respecting chirality and other stereochemical constraints. Broadly, two classes of methods are used: 

  \paragraph{Systematic or rule-based sampling:} Rotatable bonds are identified, and torsions are sampled from libraries of preferred angles; combinatorial explosion is controlled by pruning and energy filters.
  \paragraph{Stochastic sampling:} Monte Carlo or distance geometry–based methods sample conformational space randomly or quasi-randomly, often combined with local minimisation.

\subsection{OMEGA}

OMEGA (OpenEye)\cite{CH6_Hawkins2010} is a widely used systematic conformer generator for small molecules and has become a \emph{de facto} standard in many docking workflows. Its algorithm operates in three phases: first, a fragment-based initial structure assembly. In this phase, molecules are decomposed into fragments at rotatable bonds; fragment libraries derived from crystallographic data provide 3D templates. The next phase is exhaustive torsion enumeration, in which rotatable bonds are sampled at discrete torsion angles from a knowledge-based torsion library. This generates a potentially large raw conformer set. And the third phase is pruning by geometry and energy, in which conformers that are too similar (using an RMSD threshold, \textbf{Equation \ref{eq:ligRMSD}}) or too high in strain energy are discarded to yield a diverse low-energy ensemble. OMEGA\cite{CH6_Hawkins2010} has been extensively validated against high-quality crystal structures from the PDB and Cambridge Structural Database\cite{CH6_groom2016cambridge}, demonstrating high success rates in reproducing experimental ligand conformations within low RMSD thresholds.

\begin{equation}
\mathrm{RMSD} = \sqrt{\frac{1}{N} \sum_{i=1}^{N} \left\lVert \mathbf{r}_i^{(\mathrm{ref})} - \mathbf{r}_i^{(\mathrm{fit})} \right\rVert^2}
\label{eq:ligRMSD}
\end{equation}

where N is the number of matched atoms, and $\textrm{r}_\textrm{i}^{(\textrm{ref})}$ and $\textrm{r}_\textrm{i}^{(\textrm{fit})}$ are their coordinates after optimal superposition. 

\subsection{Distance geometry}

Distance geometry (DG) methods\cite{CH6_spellmeyer1997DG} (See \textbf{chapter 2} for in-depth discussion) construct 3D coordinates by embedding a matrix of pairwise distance bounds between atoms. The first step is to define lower and upper bounds $\textrm{d}_\textrm{ij}^{\textrm{min}}, \textrm{d}_\textrm{ij}^{\textrm{max}}$ for each atom pair i,j, based on covalent geometry and topological distance. Then apply triangle inequality constraints to make the bounds matrix metrically consistent. Step 3 involves sampling a distance matrix from within the bounds and performing metric matrix embedding (eigen-decomposition) to obtain coordinates. And lastly, refine the coordinates using a force-field-based minimisation.

The Experimental-Torsion Knowledge Distance Geometry (ETKDG) method \cite{CH6_Riniker2015} extends basic DG by incorporating empirical torsion-angle preferences derived from crystallographic databases and by applying additional constraints for ring systems. Implemented in RDKit, ETKDG\cite{CH6_Riniker2015} produces conformers that both reproduce experimental torsion distributions and reduce the need for extensive post-embedding minimisation. For macrocycles and highly flexible molecules, ETKDG variants\cite{CH6_wang2020improving} and specialised sampling strategies are needed to adequately explore conformational space.

\subsection{Energy minimisation and strain filtering}

Conformers produced by OMEGA, ETKDG, or other methods are often subjected to a brief energy minimisation using a molecular mechanics force field such as MMFF94\cite{CH6_halgren1996merck}, UFF\cite{CH6_rappe1992uff}, or GAFF\cite{CH6_gaff,CH6_sprenger2015general}. The typical non-bonded energy expression consists of Lennard-Jones and Coulomb terms:

\begin{equation}
E = \sum_{i<j} \left[ 4 \varepsilon_{ij} \left( \left( \frac{\sigma_{ij}}{r_{ij}} \right)^{12} - \left( \frac{\sigma_{ij}}{r_{ij}} \right)^{6} \right) + \frac{q_i q_j}{4 \pi \varepsilon_0 \varepsilon_r r_{ij}} \right],
\end{equation}

where i and j represent protein and ligand atoms respectively, $\textrm{r}_\textrm{ij}$ is the interatomic distance, $\varepsilon_\textrm{ij}$ and $\sigma_\textrm{ij}$ are Lennard--Jones parameters, $\textrm{q}_\textrm{i}$ $\textrm{q}_\textrm{j}$ are partial charges, $\varepsilon_\textrm{0}$ is the vacuum permittivity and $\varepsilon_\textrm{r}$ is the relative dielectric constants that accounts for screening effect of the medium.

Strain energy thresholds are applied to filter out highly strained conformers that would be unlikely in solution and could bias docking scores. In practice, these thresholds are empirical, typically 10-15 kcal/mol above the global minimum for a small molecule. However, they should be chosen carefully: overzealous filtering might weed out conformers relevant for binding, especially in tight or distorted pockets.

\section{Partial charge assignment and force field parameters}

\subsection{Charge models}

 The scoring functions employed in docking typically include an electrostatic term that depends on partial atomic charges. Common small-molecule charge models include empirical charge equilibration methods, such as Gasteiger–Marsili charges\cite{CH6_gasteiger1978new}, which are fast and widely available in Open Babel, RDKit, and many docking codes. There are semiempirical-based schemes, such as AM1-BCC charges\cite{CH6_jakalian2000fast}, which fit bond-charge corrections to reproduce HF/6-31G* electrostatic potentials. The RESP/ESP-fitted charges\cite{CH6_bayly1993well} are derived from higher-level QM calculations, which are common in MD force fields (GAFF/AMBER). However, these are expensive computations to handle large screening libraries. Therefore, the choice of charge model must be compatible with the docking software and, ideally, with the force field used for both the ligand and the receptor. Some protocols (e.g., Glide/SP) use internally parameterised charges and do not expose them to the user. In contrast, others (e.g., AutoDock4, Vina) require explicit charge assignment in the input (Usually, Gasteiger charges\cite{CH6_gasteiger1980iterative}).

\subsection{Parameter assignment tools}

For general organic molecules, AMBER's GAFF\cite{CH6_gaff,CH6_sprenger2015general} and CHARMM's CGenFF\cite{CH6_cgenff} are commonly used to assign force field parameters for MD, but for docking, simpler parameterisations are often sufficient. Open-source tools such as Open Babel and RDKit can assign basic atom types and MMFF94 parameters; commercial suites provide more extensive coverage of medicinal chemistry space.

\section{Software for Ligand Preparation}

\subsection{Open-source tools}

\paragraph{Open Babel:} Open Babel\cite{CH6_OBoyle2011} is a widely used open-source cheminformatics toolkit that supports interconversion between over 100 file formats, basic geometry optimisation, charge assignment, and 3D coordinate generation. It is frequently used as a main tool in ligand preparation pipelines, especially for converting between SMILES, SDF, MOL2, and PDBQT formats.

\paragraph{RDKit:} RDKit\cite{CH6_bento2020open} is a modern, open-source cheminformatics library with robust functionality for molecule sanitisation, standardisation, ETKDG-based conformer generation\cite{CH6_Riniker2015}, substructure search, and basic property calculations. Its Python API makes it an attractive choice for custom ligand preparation scripts.

\paragraph{Gypsum-DL:} Gypsum-DL\cite{CH6_Ropp2019} is an open-source program specifically designed to prepare small molecule libraries for structure-based virtual screening. Given SMILES or SDF input, Gypsum-DL:
\begin{itemize}
  \item Enumerates ionisation states (via Dimorphite-DL) over a user-defined pH range.
  \item Enumerates tautomers, R/S stereocenters, E/Z isomers, and ring conformers.
  \item Generates 3D coordinates using RDKit and filters redundant conformers.
\end{itemize}

\paragraph{Additional open-source utilities:}
\begin{itemize}
  \item Cinfony\cite{CH6_Cinfony}: a Python wrapper that unifies interfaces to Open Babel, RDKit, and CDK.
  \item Dimorphite-DL\cite{CH6_ropp2019dimorphite}: pH-dependent protonation state predictor, often used with Gypsum-DL but can be used standalone.
  \item OpenEye free academic tools: limited but valuable for generating conformers and charges.
\end{itemize}

\subsection{Commercial ligand preparation suites}

\paragraph{Schr\"odinger's LigPrep/Epik:} LigPrep, coupled with Epik for $\textrm{pK}_\textrm{a}$ prediction, is widely used in industry for systematic ligand preparation\cite{CH6_Sastry2013}. It performs:
\begin{itemize}
  \item State enumeration (protonation and tautomeric) across pH.
  \item Stereoisomer enumeration.
  \item OPLS-based\cite{CH6_jorgensen1988opls} minimisation and partial charge assignment.
  \item Export of docking-ready Maestro or MOL2 files.
\end{itemize}


\paragraph{OpenEye's QUACPAC/OMEGA:} In addition to OMEGA\cite{CH6_Hawkins2010} for conformer generation, OpenEye provides QUACPAC\cite{CH6_openeye2020quacpac} for charge state and partial charge assignment, and FILTER for removal of undesirable chemotypes. These tools are heavily used in industrial pipelines.

\paragraph{MOE, and Discovery Studio:} Both MOE (Chemical Computing Group) and Discovery Studio (BIOVIA) implement comprehensive ligand preparation workflows (salt stripping, neutralisation, protonation, stereochemistry, conformer generation) integrated with their docking engines.

\section{Verification of Ligand Structures}

\subsection{Chemistry and valence validation}

Key checks include:
\begin{itemize}
  \item Valence violations: Atoms with impossible valence (e.g., pentavalent carbon, pentavalent neutral nitrogen) or hypervalent halogens without appropriate hypervalent models.
  \item Aromaticity ring perception: Ensure aromatic rings are planar and correctly labelled; avoid inconsistent aromatic/antiaromatic assignments.
  \item Formal charge: Check for unexpected charge separation, zwitterions, or net charges inconsistent with known chemistry.
\end{itemize}

Automated sanitisation in RDKit or Open Babel can catch many such problems, but additional filters (e.g., PAINS, reactive groups) are commonly applied to drug-like libraries.

\subsection{Geometric quality}

\subsubsection{Bond lengths and angles}

Bond lengths and angles should be within reasonable ranges for the corresponding bond orders and atom types. Force field minimisation tends to correct these; however, conformations generated solely by distance geometry may show distorted geometries if poorly refined. Visual inspection of a subset of ligands and targeted checks for pathologies (e.g., extremely short non-bonded distances, highly bent triple bonds) are recommended.

\subsubsection{Ring conformations and planarity}

For aromatic and conjugated rings (e.g., benzene, heteroaromatics, amides), planarity is a useful quality criterion. Deviations from planarity can be quantified by measuring the root-mean-square deviation of ring atoms from the best-fit plane:

\begin{equation}
\mathrm{RMSD}_{\mathrm{plane}} = \sqrt{\frac{1}{N} \sum_{i=1}^{N} d_i^2},
\end{equation}

where $d_i$ is the perpendicular distance of atom i from the best-fit plane, large deviations (e.g., $>$ 0.1--0.2 \r{A} for benzene) suggest geometric artefacts. Macrocycles and flexible rings adopt non-planar conformations; however, small aromatic rings and amide linkages are expected to be close to planar in the starting conformers.

\subsection{Internal strain and unrealistic intramolecular interactions}

Excessive internal strain, such as extremely close non-bonded contacts, twisted amides, and eclipsed torsions, may artificially stabilise or destabilise certain docking poses when combined with the receptor. Minimisation with a suitable force field mitigates many such issues; however, energy-based filters can be more systematic.

\subsection{Comparison of prepared ligands to experimental structures}

For retrospective validation and method development, it is valuable to compare prepared ligand conformations and protonation states with crystallographic reference structures from PDB, PDBbind, or curated sets (DUD-E, DEKOIS 2.0). RMSD between prepared and crystallographic conformers, as well as agreement in hydrogen-bond donor/acceptor patterns and formal charges, provide quick checks\cite{CH6_Hawkins2010,CH6_WangPDBbind}. Irwin and co-workers\cite{CH6_Irwin2005} have emphasised the importance of verifying that docking scores and poses correspond to chemically reasonable structures, arguing that simple metrics like RMSD (\textbf{Equation \ref{eq:ligRMSD}}) to known structures and inspection of key interactions are essential to ensure docking is \emph{right for the right reasons}. 

\section{Databases and Datasets}

\paragraph{ZINC:}

ZINC\cite{CH6_Irwin2005} is a free database of commercially available small molecules prepared specifically for virtual screening. It provides:
\begin{itemize}
  \item Ready-to-dock 3D structures, often in multiple protonation and tautomeric states.
  \item Multiple file formats, such as SMILES, MOL2, SDF, DOCK flexibases.
  \item Annotated subsets, such as lead-like, fragment-like, and drug-like.
\end{itemize}

\paragraph{ChEMBL:}

ChEMBL\cite{CH6_Gaulton2012} is an open, manually curated bioactivity database that curates binding, functional, and ADMET data for over a million small molecules and thousands of protein targets. It is distributed with compound structures, typically as standardised SMILES or SDF, which can serve as input to ligand preparation pipelines. Efforts to curate stereochemistry, salt forms, and data quality for activity make it attractive for structure-based modelling.
For docking, ChEMBL is primarily:
\begin{itemize}
  \item A source of known actives and inactives for specific targets, useful for training, validation, and benchmarking.
  \item A source for constructing focused libraries around chemotypes of interest.
\end{itemize}

\paragraph{PubChem:}

PubChem\cite{CH6_KimPubChem}, maintained by NCBI, is a large public chemical database with substance, compound, and bioassay collections, encompassing hundreds of millions of substance records and tens of millions of unique structures. Since PubChem structures originate from many depositors, a rigorous ligand preparation and standardisation step is essential before docking. It offers:
\begin{itemize}
  \item Open access to chemical structures and associated bioactivity data.
  \item Powerful search tools, like substructure, similarity, and property filters.
  \item Links to genes, proteins, and pathways.
\end{itemize}

For docking, PubChem is often used to:
\begin{itemize}
  \item Retrieve known ligands and analogues of a target.
  \item Build custom libraries around scaffolds or pharmacophores of interest.
\end{itemize}

\paragraph{PDB and PDBbind:}

The Protein Data Bank (PDB)\cite{CH6_Berman2000} is the primary repository of macromolecular structures, including many protein–ligand complexes. These complexes define binding site geometries, water networks, and ligand poses used for docking validation and training. The PDBbind\cite{CH6_WangPDBbind} database builds on PDB by collecting experimentally measured binding affinities ($\textrm{K}_\textrm{d}$, $\textrm{K}_\textrm{i}$, and $\textrm{IC}_\textrm{50}$) for protein–ligand complexes and organising them into curated sets. PDBbind's \emph{refined set} is widely used to:

\begin{itemize}
  \item Benchmark docking and scoring functions.
  \item Train machine learning models of binding affinity.
  \item Validate ligand preparation and conformer generation against crystallographic ligands.
\end{itemize}

\paragraph{DUD-E:}

DUD-E\cite{CH6_Mysinger2012} is a benchmark library for evaluating docking and virtual screening methods. It allows assessment of enrichment metrics, such as ROC-AUC, early enrichment at 1\%, and sensitivity to ligand preparation. Mysinger \emph{et al.} showed that improved decoy selection and a more diverse set of ligands in DUD-E enhance the discriminatory power of docking benchmarks.
It comprises:
\begin{itemize}
  \item 102 protein targets with known active ligands.
  \item Property matched decoys: compounds that resemble actives in simple physicochemical properties but are topologically dissimilar and presumed inactive.
\end{itemize}

\paragraph{DEKOIS 2.0:}

DEKOIS 2.0\cite{CH6_DEKOIS2} provides challenging docking benchmark sets built from BindingDB\cite{CH6_chen2001bindingdb} bioactivity data for 81 diverse targets. Bauer \emph{et al.}\cite{CH6_DEKOIS2,CH6_IbrahimDEKOIS}  used DEKOIS 2.0 to investigate the impact of ligand and protein preparation (e.g., protonation states, conformer generation) on virtual screening performance, highlighting how different preparation protocols can significantly alter enrichment.
DEKOIS 2.0 emphasises:
\begin{itemize}
  \item Rigorous physicochemical matching of decoys to actives, including charge.
  \item Elimination of latent actives among decoys.
\end{itemize}



\paragraph{Other datasets and ML-focused docking resources:}

Beyond classical benchmarks, more recent datasets geared towards ML-based docking and scoring, such as Smiles2Dock\cite{CH6_Smiles2Dock}, provide large-scale collections of docked poses and scores for millions of ligand–protein pairs. These resources encourage development of docking surrogates and hybrid ML--physics-based approaches; however, they require careful consideration of the underlying ligand preparation methods used in dataset generation.

\section{Molecular Structure File Formats}

\subsection{SMILES}

SMILES, introduced by Weininger\cite{CH6_Weininger1988,CH6_weininger1989smiles,CH6_weininger1990smiles}, is a line notation for encoding molecular graphs as ASCII strings. It supports branching, ring closures, stereochemistry, and charges. SMILES is compact and human-readable, and is ubiquitous in cheminformatics databases and ligand preparation workflows. However, it does not contain 3D coordinates; conversion to 3D requires embedding using tools such as RDKit or Open Babel.

Examples:
\begin{itemize}
  \item Ethanol: \verb|CCO| or \verb|OCC|
  \item Benzene: \verb|c1ccccc1|
  \item Protonated amine: \verb|C[NH3+]|
\end{itemize} 

\subsection{InChI: IUPAC International Chemical Identifier}

InChI\cite{CH6_Heller2013,CH6_McNaught2006,CH6_heller2009iupac} is a non-proprietary identifier developed under IUPAC to provide a unique, canonical label for a defined chemical structure. InChI strings encode layers such as connectivity, hydrogen bonding, charge, stereochemistry, and isotopes, and InChIKey is a hashed, fixed-length representation suitable for indexing and web search. InChI is ideal for database linking and removing duplication; however, it is less convenient for substructure or similarity search than SMILES. It is not used directly as a docking input format, but is often used to harmonise ligand sets across databases (e.g., PDBbind–ChEMBL cross-mapping).

Example:
\begin{itemize}
\item InChI for ethanol is \verb|InChI=1S/C2H6O/c1-2-3/h3H,2H2,1H3|
\item InChIKey for ethanol is \verb|LFQSCWFLQFXTCV-UHFFFAOYSA-N |
\item InChI for benzene is  \verb|InChI=1S/C6H6/c1-2-4-6-5-3-1/h1-6H|
\item InChIKey for benzene is \verb|UHOVQNZJYSORNB-UHFFFAOYSA-N|
\end{itemize}

\subsection{SDF/MOL: MDL connection tables}

MDL MOL files (single molecules) and SDF (Structure Data File, multiple molecules) are widely used connection table formats. SDF allows arbitrary associated data fields (e.g., ID, SMILES, properties) and is a common interchange format for docking libraries. An SDF record for ethanol might look like:

\begin{verbatim}
Ethanol
  RDKit          2D

  3  2  0  0  0  0            999 V2000
    0.0000    0.0000    0.0000 C   0  0  0  0  0  0  0  0  0  0  0  0
    1.5400    0.0000    0.0000 C   0  0  0  0  0  0  0  0  0  0  0  0
    2.0450    1.2300    0.0000 O   0  0  0  0  0  0  0  0  0  0  0  0
  1  2  1  0
  2  3  1  0
M  END
>  <ID>
ETHANOL

>  <SMILES>
CCO

\end{verbatim}


\subsection{MOL2: Tripos MOL2 format}

MOL2 is a richer 3D format that stores atom types, charges, and substructure information. MOL2 is commonly used by docking programs such as DOCK, GOLD, and FRED, and is well supported by Open Babel and RDKit. A simplified MOL2 for ethanol might include:

\begin{verbatim}
@<TRIPOS>MOLECULE
ETHANOL
3 2 1 0 0
SMALL
NO_CHARGES

@<TRIPOS>ATOM
1 C1  0.0000  0.0000  0.0000 C.3 1 ETH 0.000
2 C2  1.5400  0.0000  0.0000 C.3 1 ETH 0.000
3 O3  2.0450  1.2300  0.0000 O.3 1 ETH 0.000

@<TRIPOS>BOND
1 1 2 1
2 2 3 1
\end{verbatim}

\subsection{PDB and PDBQT}

\begin{enumerate}
\item PDB: The PDB format was originally developed for macromolecular structures.\cite{CH6_Berman2000} Small-molecule ligands in PDB files are represented using \verb|HETATM| records. For docking, PDB is often converted to other formats (e.g., MOL2, PDBQT) that store explicit bond orders and charges more reliably than legacy PDB conventions. A minimal ligand block:

\begin{verbatim}
HETATM  1  C1  LIG A  1  10.123  12.345  5.678  1.00 20.00    C
HETATM  2  C2  LIG A  1  11.500  12.000  5.500  1.00 20.00    C
HETATM  3  O3  LIG A  1  12.000  13.200  5.800  1.00 20.00    O
\end{verbatim}
Column 2 shows atom numbers, column 3 is atom-type, column 4 is the residue name, column 5 shows the chain ID, column 5 is residue number, columns 7,8,9 are X,Y,Z coordinates respectively, column 10 is the occupancy, column 11 is temperature factor, and column 12 represents the elements.

\item PDBQT: PDBQT is an extension of PDB used by AutoDock and AutoDock Vina, adding partial charges (Q) and atom types (T) for the AutoDock force field. Bonds and rotatable torsions are encoded in separate sections. Conversion from SDF/MOL2 to PDBQT is typically done via AutoDockTools scripts or wrappers in Open Babel (\verb|obabel -i sdf -o pdbqt|).
 For example:

\begin{verbatim}
REMARK  Autodock PDBQT file
ROOT
ATOM      1  C1  LIG   1   10.123  12.345  5.678  0.00  C 
ATOM      2  C2  LIG   1   11.500  12.000  5.500  0.00  C 
BRANCH 1 2
ATOM      3  O3  LIG   1   12.000  13.200  5.800 -0.50  OA
ENDBRANCH 1 2
ENDROOT
TORSDOF 1
\end{verbatim}
\end{enumerate}

\subsection{Maestro (MAE)}

Schr\"odinger's Maestro (MAE) format and other proprietary binary/ASCII formats (e.g., SD Encrypted) store rich structural and property information. Although not human-readable, they are well supported within their respective ecosystems and often serve as the internal canonical representation during ligand preparation.

\subsection{File format considerations for docking pipelines}

Choice of file format for docking involves trade-offs:
\begin{itemize}
  \item \textbf{SMILES}: compact, ideal for initial library storage and substructure search; requires 3D embedding before docking.
  \item \textbf{SDF}: flexible and widely supported; stores multiple molecules and arbitrary properties.
  \item \textbf{MOL2}: stores atom types and charges, often preferred for docking.
  \item \textbf{PDBQT}: required by AutoDock/Vina; less general but essential in those pipelines.
\end{itemize}

\section{A Practical Ligand Preparation Workflow}

This section outlines a concrete, fully open-source workflow for preparing a library of SMILES strings for docking with a generic docking engine that accepts MOL2 or PDBQT inputs.

\paragraph{Step 1: Input and sanitisation}

\begin{enumerate}
  \item Start from a curated SMILES library, e.g., a subset of ZINC or ChEMBL filtered for basic drug-like properties.
  \item Use RDKit to parse each SMILES into a molecule object and perform sanitisation:
  \begin{itemize}
    \item Spot valence errors and discard problematic entries.
    \item Apply standardisation (functional-group normalisation, salt stripping) using predefined rules.
  \end{itemize}
\end{enumerate}

\paragraph{Step 2: Protonation and tautomer enumeration}

\begin{enumerate}
  \item For each sanitised molecule, use Gypsum-DL to enumerate protonation states and tautomers across pH 7.4 \(\pm\) 1:
  \begin{itemize}
    \item Limit the number of protonation/tautomeric states per molecule (e.g., up to 16).
    \item Reject highly unstable or chemically unreasonable states.
  \end{itemize}
  \item Alternatively, use RDKit plus an external p$K_a$ estimator or ChemAxon tools if available.
\end{enumerate}

\paragraph{Step 3: Stereochemistry and isomer enumeration}

\begin{enumerate}
  \item Within Gypsum-DL or a custom RDKit script, enumerate undefined stereocenters and E/Z isomers subject to:
  \begin{itemize}
    \item Restrictions on the maximum number of stereoisomers (e.g., up to 32).
    \item Rules to avoid redundant enumeration of centres where chirality is irrelevant (e.g., pseudoasymmetric centres).
  \end{itemize}
\end{enumerate}

\paragraph{Step 4: 3D conformer generation}

\begin{enumerate}
  \item For each enumerated structure, run RDKit's ETKDG to generate an initial ensemble of conformers (e.g., up to 20 per structure).
  \item Optionally, supplement or validate RDKit-generated conformers with OMEGA for a subset of molecules, particularly macrocycles or highly flexible chemotypes.
\end{enumerate}

\paragraph{Step 5: Energy minimisation and strain filtering}

\begin{enumerate}
  \item Minimise each conformer with MMFF94 via RDKit or Open Babel.
  \item Compute relative strain energies and discard conformers exceeding a threshold (e.g., 10 kcal/mol above the minimum).
  \item Optionally cluster conformers based on heavy-atom RMSD to ensure diversity.
\end{enumerate}

\paragraph{Step 6: Charge assignment}

\begin{enumerate}
  \item If targeting AutoDock/Vina, assign Gasteiger partial charges using Open Babel or AutoDockTools.
  \item For other docking engines, use MMFF94 or AM1-BCC charges if supported, or the engine's internal charge model.
\end{enumerate}

\paragraph{Step 7: Export to docking format}

\begin{enumerate}
  \item Export per-structure conformers to SDF or MOL2 using Open Babel.
  \item For AutoDock/Vina, convert MOL2/SDF to PDBQT and define rotatable bonds.
\end{enumerate}

\paragraph{Step 8: Quality Check}

\begin{enumerate}
  \item Randomly sample a few hundred ligands and visually inspect:
  \begin{itemize}
    \item Protonation states and tautomers (do they match expected chemistry?).
    \item Ring conformations and planarity of aromatics and amides.
  \end{itemize}
  \item Run automated checks:
  \begin{itemize}
    \item Validate valence, bond orders, and charges.
    \item Detect severe steric clashes or atypical geometries.
  \end{itemize}
\end{enumerate}

\section{Best Practices}

Robust small-molecule preparation is a prerequisite for meaningful docking-based virtual screening. Key recommendations are:

\begin{itemize}
  \item Treat ligand preparation as a modelling step, not a trivial pre-processing task. Empirical studies show that preparation parameters can substantially affect enrichment. 
  \item Use chemically informed protonation and tautomer states. Incorporate p$K_a$ prediction and tautomer enumeration; avoid assuming a single arbitrary state for multi-protic systems.
  \item Enumerate relevant stereochemistry. Do not ignore undefined stereocenters if stereochemistry is expected to affect binding.
  \item Employ validated conformer generators. Use tools such as OMEGA and RDKit/ETKDG, which have been benchmarked against crystallographic data.
  \item Filter out chemically implausible and highly strained conformers. Apply energetic and geometric criteria to remove outliers.
  \item Leverage curated databases. Build libraries from ZINC, ChEMBL, PubChem, PDBbind, and established benchmarks (DUD-E, DEKOIS 2.0) to maximise data quality and facilitate comparison with literature.
  \item Verify. Combine automatic checks with manual inspection of representative ligands and docked poses to ensure that the modelling remains chemically reasonable.
\end{itemize}

\bibliographystyle{unsrtnat}
\bibliography{chapter6}

\end{document}
